\chapter{Introducción}

El problema abordado en este laboratorio corresponde a la descripción comparativa de desigualdad territorial entre comunas de la Provincia de Santiago, usando tres dimensiones sustantivas: disponibilidad de áreas verdes, pobreza por ingresos y capacidad municipal. La unidad de análisis observada en el proyecto es comunal, por lo que cada registro final representa una comuna y resume variables provenientes de distintas instituciones públicas.

La construcción de un proceso ETL resulta necesaria porque las fuentes disponibles no llegan en un formato analítico homogéneo. En el repositorio se observan diferencias de formato de archivo, estructuras de encabezado, escalas de medición y variantes ortográficas en \texttt{nombre\_comuna}. En particular, dos fuentes SINIM fueron publicadas como archivos con extensión \texttt{.xls} pero estructura SpreadsheetML/XML 2003, mientras que las otras dos corresponden a libros \texttt{.xlsx}. Adicionalmente, los nombres comunales cambian entre mayúsculas, tildes y escritura capitalizada, lo que vuelve insuficiente cualquier integración basada en nombre crudo.

El enfoque ETL implementado en el proyecto separa explícitamente las etapas de perfilado, homologación, extracción a staging, transformación por fuente, integración final, validación estructural, carga a SQLite y análisis exploratorio. Esta secuencia permite preservar trazabilidad entre insumos y resultados, registrar decisiones metodológicas y documentar de manera reproducible cómo una combinación de fuentes heterogéneas se transforma en un dataset comunal consistente.

El proyecto integra múltiples fuentes públicas porque ninguna de ellas contiene por sí sola la unidad analítica final requerida. La población comunal provee el denominador para indicadores por habitante, la pobreza por ingresos aporta una dimensión socioeconómica, las áreas verdes entregan una dimensión territorial y los ingresos propios permanentes aportan una aproximación a capacidad municipal. El valor del producto final no reside en inferir relaciones causales, sino en dejar una base integrada, consistente y consultable para describir desigualdades comunales observadas.
